%==========================================================================================
\section{Older and superseded tools}\label{sec:legacy}
%==========================================================================================

These modules predate the current analysis chain. They are kept in the repository for
reference and for old data, and documented here with the problems that stop them from running
unchanged on a current installation (Python~3.12, NumPy~2).

%------------------------------------------------------------------------------------------
\subsection{\texttt{fluo\_map\_fit}: the first fluorescence map fitter}
%------------------------------------------------------------------------------------------

Fits every spectrum of a BLISS map independently with two ``coupled double Gaussians'' (Ga
and In), a constant background and \fn{curve_fit}, in parallel over the pixels, and plots the
two area maps. Each element is a K$\alpha$ peak plus a K$\beta$ peak at a fixed channel
distance and a fixed amplitude ratio, sharing one width:
\begin{equation}
  g(x) = A\,\e^{-(x-c)^2/2\sigma^2} + rA\,\e^{-(x-c-\Delta)^2/2\sigma^2},
  \qquad \text{area} = A\sigma\sqrt{2\pi}\,(1+r).
\end{equation}

\begin{issue}
Superseded by \fn{xrf_fit} (Section~\ref{sec:fit}), which replaces every hard-coded number
here by tabulated physics. The channel offsets and ratios are hard-coded for a 10\,eV/channel
calibration (Ga: $\Delta = 101$ channels, $r = 0.13$; the In values are placeholders), as are
the fitted channel range (800--1330), the starting centres and their bounds, and the input
file in the \fn{__main__} block. The spectra are not normalised (monitor, live time), the
widths do not follow the detector resolution law, and K$\beta$/K$\alpha$ is not computed for
the excitation. The module docstring calls it \fn{constrained_xrf_fitter.py}.
\end{issue}

\pycode{fluo_map_fit}{_header}

\api[\texttt{coupled\_double\_gaussian}, \texttt{spectrum\_model}]{\texttt{coupled\_double\_gaussian(x, amp\_alpha, cen\_alpha, sigma, delta, ratio)}, \texttt{spectrum\_model(x, bkg, amp\_ga, cen\_ga, sig\_ga, amp\_in, cen\_in, sig\_in)}}
The doublet above, and the full model: constant background plus the Ga and In doublets.
\pycode{fluo_map_fit}{coupled_double_gaussian}
\pycode{fluo_map_fit}{spectrum_model}

\api[\texttt{fit\_single\_spectrum}, \texttt{process\_and\_plot}]{\texttt{fit\_single\_spectrum(args)}, \texttt{process\_and\_plot(h5\_filepath, dataset\_path, scan\_shape)}}
\fn{fit_single_spectrum} fits one pixel, \fn{args = (index, channels, counts)}, and returns
\fn{(index, Ga area, In area)} (both \fn{None} if the fit fails). \fn{process_and_plot} loads
the MCA dataset, slices the channel range, distributes the pixels over a
\fn{ProcessPoolExecutor} (one core left free, chunks of 100), reshapes the areas to
\fn{scan_shape} and plots the two maps.
\pycode{fluo_map_fit}{fit_single_spectrum}
\pycode{fluo_map_fit}{process_and_plot}
\pycode{fluo_map_fit}{_tail}

%------------------------------------------------------------------------------------------
\subsection{\texttt{spec\_reader}: SPEC data files}
%------------------------------------------------------------------------------------------

SPEC, the control program used on many synchrotron beamlines before BLISS, writes plain-text
data files: a file header (\texttt{\#F} file name, \texttt{\#E} epoch, \texttt{\#D} date, \texttt{\#O} motor
names), then for each scan a header (\texttt{\#S} number and command, \texttt{\#D} date, \texttt{\#T} count
time, \texttt{\#G} geometry, \texttt{\#Q} $hkl$, \texttt{\#P} motor positions, \texttt{\#N} number of columns,
\texttt{\#L} column labels, \texttt{\#C} comments, \texttt{\#@} MCA and user lines), followed by the data
columns. \fn{SpecFile} reads the file header and indexes the byte offset of every scan;
\fn{Scan} reads one scan---or concatenates several similar scans---and exposes every counter
as an attribute.

\begin{issue}
\begin{itemize}[nosep,leftmargin=1em]
\item Files are opened with mode \fn{'rU'}, which was removed in Python~3.11: on the current
  environment \fn{SpecFile} raises \fn{ValueError: invalid mode: 'rU'} (not caught, since
  only \fn{IOError} is).
  Mode \fn{'r'} gives the same universal-newline behaviour in Python~3.
\item \fn{__detcalib__} uses \fn{np.float}, removed in NumPy~1.24 (only reached for
  \texttt{\#UDETCALIB} lines).
\item The XIA handlers are unreachable: a line \texttt{\#@XIAFILE path} gives the key
  \fn{@XIAFILE}, which is not in the dispatch table, and a line \texttt{\#@ XIAFILE path} reaches
  \fn{__special__} but splits into an empty first word. Both end as ``unprocessed lines''
  (checked on the current code), so \fn{xianame}, \fn{xiacalib} and \fn{xiaroi} are never
  set.
\end{itemize}
\end{issue}

\api[\texttt{SpecFile}]{\texttt{SpecFile(spec\_file, verbose=False)}}
Reads the file header up to the first \texttt{\#S} line, dispatching each line to a handler by its
key (\fn{__param__} maps the key letters to methods), and records in \fn{scan_dict} the byte
position of every \texttt{\#S} line, so that a scan can later be read with a single \fn{seek}.
The key is the text after \texttt{\#}, with trailing digits removed except for user keys starting
with \fn{U} (so \texttt{\#O0}, \texttt{\#O1} both go to the \fn{O} handler).
\pycode{spec_reader}{_header}
\pycode{spec_reader}{SpecFile}
\pycode{spec_reader}{SpecFile.__param__}
\pycode{spec_reader}{SpecFile.__readSpecLine__}
\pycode{spec_reader}{SpecFile.__init__}

\api{Line handlers}
One method per key, each parsing its line into attributes:
\fn{__scanline__} (\texttt{\#S}: \fn{number}, \fn{type}, \fn{args}, \fn{command}),
\fn{__dating__} (\texttt{\#D}), \fn{__counting__} (\texttt{\#T}: \fn{ct}, \fn{ct_units}),
\fn{__configurating__} (\texttt{\#G}), \fn{__hkl__} (\texttt{\#Q}: \fn{Qo}),
\fn{__motorslabeling__} and \fn{__motorslabelingnospace__} (\texttt{\#O}, \texttt{\#o}),
\fn{__positioning__} (\texttt{\#P}), \fn{__speccol__} (\texttt{\#N}: \fn{N}),
\fn{__counterslabeling__} (\texttt{\#L}: \fn{counters}), \fn{__allcounterslabeling__} and
\fn{__allcounterslabelingnospace__} (\texttt{\#J}, \texttt{\#j}), \fn{__specfilenaming__} (\texttt{\#F}),
\fn{__commenting__} (\texttt{\#C}), \fn{__initialepoch__} (\texttt{\#E}), \fn{__marccdpath__}
(\texttt{\#M}), \fn{__limampx4path__} (\texttt{\#ULIMA\_mpx4}), \fn{__detcalib__} (\texttt{\#UDETCALIB}:
detector centre, pixels per degree, distances), \fn{__special__} (\texttt{\#@}: dispatches on the
following word), \fn{__xiafilenaming__}, \fn{__xiacalibrating__}, \fn{__xiaroidefining__}
(XIA MCA file, calibration and ROIs).
\pycode{spec_reader}{SpecFile.__scanline__}
\pycode{spec_reader}{SpecFile.__dating__}
\pycode{spec_reader}{SpecFile.__counting__}
\pycode{spec_reader}{SpecFile.__configurating__}
\pycode{spec_reader}{SpecFile.__hkl__}
\pycode{spec_reader}{SpecFile.__motorslabeling__}
\pycode{spec_reader}{SpecFile.__motorslabelingnospace__}
\pycode{spec_reader}{SpecFile.__positioning__}
\pycode{spec_reader}{SpecFile.__speccol__}
\pycode{spec_reader}{SpecFile.__counterslabeling__}
\pycode{spec_reader}{SpecFile.__allcounterslabeling__}
\pycode{spec_reader}{SpecFile.__allcounterslabelingnospace__}
\pycode{spec_reader}{SpecFile.__specfilenaming__}
\pycode{spec_reader}{SpecFile.__commenting__}
\pycode{spec_reader}{SpecFile.__initialepoch__}
\pycode{spec_reader}{SpecFile.__marccdpath__}
\pycode{spec_reader}{SpecFile.__limampx4path__}
\pycode{spec_reader}{SpecFile.__detcalib__}
\pycode{spec_reader}{SpecFile.__special__}
\pycode{spec_reader}{SpecFile.__xiafilenaming__}
\pycode{spec_reader}{SpecFile.__xiacalibrating__}
\pycode{spec_reader}{SpecFile.__xiaroidefining__}

\api[\texttt{Scan}]{\texttt{Scan(spec\_file, scan\_numbers, verbose=False)}}
Reads one scan, or several concatenated, from a \fn{SpecFile} (or a file name).
\begin{args}
\item[spec\_file] a \fn{SpecFile}, or a path (then indexed first).
\item[scan\_numbers] one scan number, or a sequence of them.
\item[verbose] print the lines that are read or not processed.
\end{args}
\blk{Method} The first scan's header is parsed with the handlers; data lines are read until a
blank line or a \texttt{\#} line other than \texttt{\#C}. Each further scan is appended if it is
``similar'': same scan type, same first argument (the scanned motor) and, for commands with
more than five arguments, the same fourth argument (for an \fn{ascan}, the number of
intervals); otherwise a message is printed and the scan is skipped. Every column becomes an
attribute named after its counter label, \fn{motors} maps the motor names to their start
positions (using the \texttt{\#o} names when motor names contain spaces), and with an \fn{Epoch}
column the start and end times, duration and time per point are set. A mismatch between
\texttt{\#N} and the number of labels is reported.
\pycode{spec_reader}{Scan}
\pycode{spec_reader}{Scan.__init__}
\pycode{spec_reader}{_tail}

%------------------------------------------------------------------------------------------
\subsection{\texttt{xray\_tools}: wavelengths, Bragg angles and flipping ratios}
%------------------------------------------------------------------------------------------

Small tools from the SPEC era, including the analysis of X-ray phase retarder (XPR) scans: a
diamond crystal rocked near a Bragg reflection turns the linear polarisation of the beam into
circular polarisation of opposite helicities on either side of the reflection
\citep{Hirano1991}, and the \emph{flipping ratio} of a signal measured at symmetric offsets is
the helicity asymmetry (e.g. an XMCD signal).

\begin{issue}
\begin{itemize}[nosep,leftmargin=1em]
\item \fn{xprplot_th} computes \fn{midpoint = n/2}, a float in Python~3, then slices with it:
  \fn{TypeError}. Integer division (\fn{//}) is needed.
\item Only \fn{matplotlib} is imported, but \fn{matplotlib.pyplot} is used: it fails unless
  \fn{pyplot} was imported elsewhere first.
\item The module does not import on the current environment: \fn{from numpy.linalg.linalg
  import *} raises \fn{ModuleNotFoundError} with NumPy~2.5 (the submodule was made private in
  NumPy~2.0; \fn{from numpy.linalg import norm} is the replacement). \fn{from numpy import *}
  also puts the whole NumPy namespace into the module.
\item Scans are read with \fn{spec_reader}, which cannot open files on Python~$\ge$~3.11.
\end{itemize}
\end{issue}

\pycode{xray_tools}{_header}

\api[\texttt{E2L}, \texttt{L2E}]{\texttt{E2L(energy)}, \texttt{L2E(wavelength)}}
$\lambda = hc/E$ and $E = hc/\lambda$, in metres and electronvolts, with
$h = 4.135667516\times10^{-15}$\,eV\,s and $c = 2.99792458\times10^8$\,m/s
(CODATA 2010, \citealp{Mohr2012}): $hc = 1.2398\times10^{-6}$\,eV\,m.
\pycode{xray_tools}{E2L}
\pycode{xray_tools}{L2E}

\api[\texttt{tth}]{\texttt{tth(H, E, a)}}
Bragg angle $2\theta$ (degrees) of the reflection $\mathbf H = (h,k,l)$ of a cubic crystal of
lattice parameter $a$ (m) at energy $E$ (eV): $d = a/|\mathbf H|$ and
$2\theta = 2\arcsin(\lambda/2d)$ \citep{Bragg1913}. Example: (0,0,10) at 10\,871\,eV with
$a = 10.015$\,\AA\ gives 69.42\textdegree.
\pycode{xray_tools}{tth}

\api[\texttt{uncertainty\_fr}]{\texttt{uncertainty\_fr(p, n, po, no)}}
Flipping ratio and its uncertainty from the detector counts $p$, $n$ and monitor counts
$p_o$, $n_o$ for the two helicities:
\begin{equation}
  \mathrm{FR} = \frac{P-N}{P+N}, \quad P = \frac{p}{p_o},\ N = \frac{n}{n_o}; \qquad
  \delta\mathrm{FR} = \sum_{x\in\{p,n,p_o,n_o\}}\left|\frac{\partial\,\mathrm{FR}}{\partial x}\right|
  \sqrt x ,
\end{equation}
with, for instance, $\partial\mathrm{FR}/\partial p = 2N/[p_o(P+N)^2]$. Returns
\fn{(dfr, fr)}.
\begin{hyp}
Poisson statistics for all four counts. The contributions are added linearly, which is an
upper bound; independent errors would add in quadrature \citep{Bevington2003}, giving a value
up to two times smaller.
\end{hyp}
\pycode{xray_tools}{uncertainty_fr}

\api[\texttt{xprplot\_th}]{\texttt{xprplot\_th(spec\_file, scan\_number, detector, monitor, centered=True, do\_plot=True, save=False, auto\_close=False)}}
Processes an XPR angle scan measured symmetrically about the centre of the phase-retarder
Bragg peak: the second half of the scan gives one helicity ($p$, $p_o$) and the first half,
reversed, the other ($n$, $n_o$), at the same offsets; the centre point is skipped for an odd
number of points. Computes the flipping ratio at each offset, optionally plots and saves it
(PDF, EMF, PNG). Returns \fn{(offsets, fr, dfr)}. Non-centred scans are not implemented.
\pycode{xray_tools}{xprplot_th}

\api[\texttt{XprScan}, \texttt{XprScan.fr}]{\texttt{XprScan}, \texttt{XprScan.fr(detector, monitor, ...)}}
A \fn{spec_reader.Scan} with a method that runs \fn{xprplot_th} on itself and stores the
offsets, flipping ratio and error as attributes (\fn{xpr_offset}, \fn{<detector>_fr},
\fn{<detector>_dfr}).
\pycode{xray_tools}{XprScan}
\pycode{xray_tools}{XprScan.fr}

%------------------------------------------------------------------------------------------
\subsection{\texttt{LiveFourierTransform}: webcam Fourier transform}
%------------------------------------------------------------------------------------------

A demonstration script (2011, from a Matlab original by B.~R. Pauw) that grabs webcam frames,
computes the two-dimensional Fourier transform of a region of interest and displays, side by
side, the real-space image, the logarithm of the intensity $|F|^2$, of the modulus $|F|$, and
the phase of the transform, each scaled to 8 bits. It illustrates that a far-field
(Fraunhofer) diffraction pattern is the squared modulus of the Fourier transform of the object
\citep{Goodman2005}, and that the phase---lost in a diffraction measurement---carries much of
the information. The transform uses the FFT \citep{Cooley1965} with the zero frequency shifted
to the centre; the optional ``kill centre lines'' replaces the central row and column (which
hold the strong contributions of the image borders) by their neighbours to increase the
dynamic range.

\begin{issue}
Python~2 code (\fn{print} statement, integer division \fn{h/2} for indexing) using the
OpenCV~1.x Python bindings (\fn{opencv}, \fn{highgui}) and the old \fn{Image} module of PIL;
\fn{numpy} and its functions (\fn{angle}, \fn{log10}, \fn{ones}, \fn{uint8}, \fn{transpose},
\fn{append}) are used without being imported. It does not run on a current installation and
is kept as a historical record. Because it is not valid Python~3, the listing shows the whole
file.
\end{issue}

\pycode{LiveFourierTransform}{_file}
